\centerline{\bf \Huge Комплы. Практика}

\begin{flushright}
        \scriptsize{}
\end{flushright}

\problem{1}
В результате поворота на $90^{\circ}$ вокруг точки $O$ отрезок $A B$ перешел в $A_1 B_1$. Докажите, что медиана $O M$ треугольника $O A B_1$ перпендикулярна прямой $A_1 B$.

\problem{2}
Дан вписанный четырёхугольник $A B C D$. Обозначим через $H_a, H_b, H_c, H_d$ ортоцентры треугольников $B C D, A C D, A B D, A B C$ соответственно. Докажите, что прямые $A H_a, B H_b, C H_c, D H_d$ пересекаются в одной точке.

\problem{3}
Прямая $t$ касается описанной окружности треугольника $A B C$ в точке $B . K$ проекция ортоцентра $A B C$ на $t, L$ середина $A C$. Докажите, что треугольник $B K L$ равнобедренный.

\problem{4}
Докажите, что середины трех отрезков, соединяющих проекции произвольной точки плоскости на пары противоположных сторон или диагоналей вписанного в окружность четырехугольника, лежат на одной прямой.

\problem{5} 
$\bigl[\textbf{Отбор Эйлера}\bigr]$
Внутри трапеции $A B C D(B C \| A D)$, где $A D=2 B C$, взята точка $F$, для которой $A B=F B$. Точка $M$ - середина отрезка $F D$. Докажите, что $C M \perp F A$.

\problem{6}
$\bigl[\textbf{Регион Эйлера}\bigr]$
Точка $N$ - середина стороны $B C$ треугольника $A B C$, в котором $\angle A C B=60^{\circ}$. Точка $M$ на стороне $A C$ такова, что $A M=B N$. Точка $K$ - середина отрезка $BM$. Докажите, что $A K=K C$.

\problem{7}
$\bigl[\textbf{ММО, 10.3}\bigr]$
Точка $O$ - центр описанной окружности треугольника $A B C, A H$ - его высота. Точка $P$ - основание перпендикуляра, опущенного из $A$ на прямую $C O$. Докажите, что прямая $H P$ проходит через середину отрезка $AB$.

\problem{8}
$\bigl[\textbf{Сложный тургор, младшие, 4 задача}\bigr]$
На сторонах $B C$ и $C D$ ромба $A B C D$ отмечены точки $P$ и $Q$ соответственно так, что $B P=C Q$. Докажите, что точка пересечения медиан треугольника $A P Q$ лежит на диагонали $B D$ ромба.